\documentclass[10pt,a4paper]{amsart}
\usepackage[margin=19mm]{geometry}
\usepackage{amsmath,amssymb,mathtools}
\usepackage[scr=rsfs,cal=euler]{mathalfa}
\usepackage{xfrac}
\usepackage[unicode,hidelinks,bookmarks=false]{hyperref}
\newcommand{\ie}{i.e.\ }
\newcommand{\cf}{cf.\ }
\newcommand{\resp}{respectively\ }
\newcommand{\Subsection}{\subsection}
\providecommand{\nobreakdash}{\nobreak\hbox{-}}
\setlength{\emergencystretch}{2em}
\multlinegap0pt
\usepackage{amssymb}
\usepackage{mathtools}
\usepackage{enumerate}
\usepackage{stackengine}
\usepackage{tikz}
\usepackage{xcolor}
\newtheorem{theorem}{Theorem}
\newcommand{\N}{\mathbb{N}}
\newcommand{\R}{\mathbb{R}}
\title{Growth beyond exponent $3/2$ for convexity and iterated sum sets}
\author{Michalis Kokkinos, Oliver Roche-Newton}
\date{Source: arXiv:2609.15265v1 (CC BY 4.0)}
\begin{document}
\maketitle

\noindent\emph{Source and licence.} Growth beyond exponent $3/2$ for convexity and iterated sum sets. Version arXiv:2609.15265v1 (CC BY 4.0), distributed by arXiv under the Creative Commons Attribution 4.0 International licence, http://creativecommons.org/licenses/by/4.0/, as recorded in the arXiv licence metadata for this exact version and shown on the abstract page https://arxiv.org/abs/2609.15265v1. Full licence notice: \texttt{licenses/arxiv-2609.15265-CC-BY-4.0.txt}.\par\smallskip

\noindent\textit{Editorial scope.} Counted unit: the article's theorem generalpoly (source0.tex lines 169--180). The article's complete original proof of it is delivered with it (source0.tex lines 181--200, one \texttt{proof} environment of 20 lines). No mathematics is written, repaired or re-derived: the statement and the proof are the article's own lines, in the article's own notation, and no retained line is substituted or re-flowed.  Retained uncounted as the source-local context the proof needs: lines 136--149; lines 157--162.  Nothing was omitted from any retained span, so the package carries no omission record.  No retained line contains a citation, so the wrapper carries no bibliography and no bibliography entry.  No retained line contains a figure environment, an image-inclusion command or drawing code, so no image is delivered and the package declares no asset dependency.  Editorial formatting only, in addition: an AMS \texttt{amsart} preamble that restates the article's own declarations and macros used by the retained lines, namely the environments \texttt{theorem} on separate counters, together with the standard packages those lines need.  The retained environments print in this document as Theorem 1 in order of appearance, whereas the article prints the same items with the numbering of its own full document.  The complete native source of the article as distributed in the arXiv e-print for this exact version is bundled unchanged as \texttt{sources/210404/source0.tex}.
\begin{theorem}\label{core}
    Let $I\subseteq\R$ be an interval of real numbers and let $m\in\N$. Suppose
that $A\subseteq I$ is a finite set of reals and that $f : I \rightarrow\R$ is a strictly convex or concave function. For $d\in(0,\infty)$, let $I_d=\{x\in I:x+d\in I\}$ and define the function $f_d:I_d\rightarrow\R$ by the formula $f_d(x) = f(x+d)-f(x)$. Suppose that, for all $d \in (0, \infty)$, the first four derivatives of $f_d^{-1}:R_{f_d}\rightarrow I_d$ have a combined total of at most $M$ zeroes. Then 
\[
    \max \{|16A|, |13f(A)| \} \gg_M |A|^{\frac{3}{2} +\frac{1}{162}}.
    \]
%\begin{equation*}
 %   |8A-7A|^{16}|5f(A)-4f(A)|^{11}\gg_m\frac{|A|^{41}}{(\textup{log}|A|)^{77}}
%\end{equation*}
%In particular,
%\begin{equation*}
 %   \text{max}\{|8A-7A|,|5f(A)-4f(A)|\}\gtrsim_m|A|^{\frac{3}{2}+\frac{1}{54}}.
%    \end{equation*}
\end{theorem}

\begin{align} 
\big(f_d^{-1}\big)^{(1)}(y)&=\frac{1}{f_d^{(1)}(x)}, \label{d1} \\
    \big(f_d^{-1}\big)^{(2)}(y)&=\frac{-f_d^{(2)}(x)}{\big(f_d^{(1)}(x)\big)^3},  \label{d2}\\ 
    \big(f_d^{-1}\big)^{(3)}(y)&=\frac{-f_d^{(3)}(x)f_d^{(1)}(x)+3\big(f_d^{(2)}(x)\big)^2}{\big(f_d^{(1)}(x)\big)^5},\text{ and} \label{d3} \\
    \big(f_d^{-1}\big)^{(4)}(y)&=\frac{10f_d^{(1)}(x)f_d^{(2)}(x)f_d^{(3)}(x)-\big(f_d^{(1)}(x)\big)^2f_d^{(4)}(x)-15\big(f_d^{(2)}(x)\big)^3}{\big(f_d^{(1)}(x)\big)^7}. \label{d4}
    \end{align}

\begin{theorem}\label{generalpoly}
    Let $A$ be a finite set of reals. Let $f:\R\rightarrow\R$ be a polynomial function given by $f(x)= a_mx^m+a_{m-1}
x^{m-1}+\dotsc+a_1x+a_0\in\R[x]$ of degree $m\ge3$, and with $a_m\neq0$. Then, 
%there exists interval $I\subseteq\R$ for which $h:=f\big|_I:I\rightarrow\R$ is a strictly convex or concave function and $A'\subseteq A\cap I$, with $|A'|\ge|A|/n$, such that
%\begin{equation*}
 %   |8A-7A|^{16}|5f(A)-4f(A)|^{11}\gg_n\frac{|A|^{41}}{(\textup{log}|A|)^{77}}
%\end{equation*}
%In particular,
  \[
    \max \{|16A|, |13f(A)| \} \gg_m |A|^{\frac{3}{2} +\frac{1}{162}}.
    \]
\end{theorem}

\begin{proof}
    We start by observing that the function $f_d$ defined as indicated above is a polynomial function of degree $m-1$. The zeroes of $(f_d^{-1})^{(i)}$ occur precisely when the numerators of the expressions \eqref{d2}, \eqref{d3} and \eqref{d4} are zero. These numerators are polynomials with degree at most $m^3$, and hence they have at most $m^3$ zeroes unless they are identically zero. Hence, the technical condition on $f$ of Theorem \ref{core} is satisfied if we prove that none of the numerators of the derivatives in \eqref{d2}, \eqref{d3}, \eqref{d4} can be identically zero. 
    
    We consider the highest order term in the numerator of each of these three derivatives. For \eqref{d2}, the numerator is a polynomial of degree $m-3 \geq 0$, as required (note that in the case $m=3$ this is a non-zero constant polynomial). 

    The highest order term for the numerator in \eqref{d3} is 
     %Note that $f_d(x)=a_nndx^{n-1}\,+$ lower order terms. Indeed, the highest order term of the numerator of the second derivative $-a_nn(n-1)(n-2)dx^{n-3}$, where $a_n\neq0$, $d>0$ and $n\ge3$ by assumption, is not zero for all $x\in I.$ For the third derivative, we consider the highest order term of $f_d^{(3)}(x)f_d^{(1)}(x)+3\big(f_d^{(2)}(x)\big)^2$, which is
     \begin{equation*}
        a_m^2m^2(m-1)^2(m-2)d^2\Big[-(m-3)+3(m-2)\Big]x^{2m-6}.
    \end{equation*}
    This vanishes for all $x\in \R$ precisely when $m-3=3(m-2)$, that is when $m=3/2$, which cannot occur. Lastly, we proceed to check for the non-vanishing of the numerator of \eqref{d4}. Here, the leading coefficient is \begin{equation*}
        \begin{split}
            a_m^3m^3(m-1)^3(m-2)d^3\Big[10(m-2)(m-3)-(m-3)(m-4)-15(m-2)^2\Big].
        \end{split}
    \end{equation*}
    We deduce that the leading coefficient is not zero, since this requires $6m^2-17m+12=0,$ that is $m\in\{3/2,4/3\}$. We have checked that $f$ satisfies the derivative conditions of Theorem \ref{core}. 
    % If $m=3,$ then the highest order term of the numerator of \eqref{d4} is the constant term $-15\big(f_d^{(2)}(x)\big)^3=-15(6a_3d)^3\neq0.$
    
    We note that we cannot yet apply Theorem \ref{core} directly since the function $f$ as defined above is not strictly convex or concave. However, since $f''$ has degree $m-2$, the real line can be partitioned into $t$ intervals $C_i\subseteq\R$, for some $t\le m-1$, on which $f$ is strictly convex or concave. Then $|A|=\sum_{i=1}^t|C_i\cap A|$ and, by the Pigeonhole Principle, there exist $j\in\{1,\dotsc,t\}$ for which $|C_j\cap A|\ge|A|/t\gg_m|A|.$ We can now apply Theorem \ref{core} with the set of real numbers chosen to be the large intersection $C_j\cap A$ and with the strictly convex or concave function $f\big|_{C_j}$, since the relevant sumsets of the subset are contained in those of $A.$
\end{proof}

\end{document}
